% !TEX root = ../../main_power.tex
% ---------------------------------------------------------------------
%  Leistung und Energie -- die physikalische Grundlage (Teil 3)
% ---------------------------------------------------------------------
\section{Leistung und Energie}
\label{sec:pwr}

\begin{formel}[Die zwei Grundgleichungen]
\[ P = U \cdot I \qquad
   E = \int_0^T P(t)\,dt = \int_0^T I(t) \cdot V(t)\,dt \]
\merkzeile Leistung ist ein \emph{Momentanwert}, Energie die
\emph{Fläche darunter}. Ein Messgerät, das nur den Mittelwert zeigt,
sagt nichts über Spitzen -- und umgekehrt.
\end{formel}

\SW{„Low Power“ ist nicht „Low Energy“} Die beiden Grössen beantworten
verschiedene Fragen, und eine Massnahme kann die eine verbessern und die
andere verschlechtern.
\begin{ztab}{@{}lX@{}}
Leistung & bestimmt das \emph{Design}: Speisung, Spannungsregler,
  Leiterquerschnitte, Kühlung \\
Energie & bestimmt den \emph{Betrieb}: Batteriekapazität, Energiekosten
  (Solar, Weltraum), Kühlung (Platz und Kosten), Lebensdauer
  (temperaturabhängig) \\
\end{ztab}
\merkzeile Darum lohnt sich „schnell rechnen und lange schlafen“: die
Spitzenleistung steigt, die Energie sinkt trotzdem.

\subsection{Statischer Verbrauch}
\label{sec:statisch}
Fliesst, \emph{ohne} dass geschaltet wird.
\begin{itemize}
\item Leck-/Kanalstrom von Source nach Drain
\item Leckstrom durch die parasitäre Diode
\item endlicher Widerstand des Halbleiters
\end{itemize}
\begin{falle}[Die Zwickmühle des statischen Verbrauchs]
Gegenmittel ist die \T{angelegte Spannung senken}. Aber: reduzierte
Spannung \ra längere Schaltzeiten. Will man die Schaltzeiten zurück,
muss man die \T{Schaltspannung} (Schwellspannung) senken -- und dann
steigt der Leckstrom \T{exponentiell}. Statischer Verbrauch lässt sich
nicht wegoptimieren, nur verschieben.
\end{falle}

\subsection{Dynamischer Verbrauch}
\label{sec:dynamisch}
Fliesst nur \emph{beim} Schalten, hat drei Anteile:
\begin{ztab}{@{}lX@{}}
$I_{SW}$ & Laden und Entladen der Lastkapazität am Ausgang -- der
  Hauptanteil \\
$I_{SC}$ & Kurzschlussstrom: für kurze Zeit leiten \emph{beide}
  Transistoren, weil die Schaltflanke nicht unendlich steil sein kann \\
$I_{lk}$ & Leckstrom (siehe oben) \\
\end{ztab}
\begin{formel}[Warum Kapazität und Takt zählen]
\[ I = \frac{dQ}{dt} = C \cdot \frac{dV}{dt} \]
Der Strom ist proportional zur Kapazität $C$ \emph{und} zur
Schaltgeschwindigkeit $dV/dt$. Bei Millionen von Transistoren, die alle
geladen und entladen werden, summiert sich das.
\end{formel}

\SW{Trend} Kleinere Strukturen \ra kleinere Kapazitäten \ra dynamischer
Verbrauch sinkt, \emph{statischer steigt}. Höhere Taktrate \ra
dynamischer Verbrauch steigt, aber der Einfluss der Kapazität ist
grösser. Das eigentliche Problem der Miniaturisierung sind die
Leckströme: irgendwann kreuzen sich die beiden Kurven.

\subsection{Die Gleichung, aus der alle Massnahmen folgen}
\label{form:ptotal}
\begin{formel}[Gesamtverbrauch]
\[ P_{total} = \underbrace{\alpha\,C_{load}\,V_{dd}^{2}\,f}_{\text{dynamisch}}
 + \underbrace{t_{SC}\,V_{dd}\,I_{peak}}_{\text{Kurzschluss}}
 + \underbrace{V_{dd}\,I_{leak}}_{\text{statisch}} \]
\end{formel}
Jeder Faktor ist eine Stellschraube -- die ganze Vorlesung besteht
daraus, sie einzeln anzuziehen:
\begin{ztab}{@{}lX@{}}
$V_{dd}$ & Speisespannung senken, abschalten, interne Spannung nutzen.
  Wirkt \T{quadratisch} im dynamischen \emph{und} linear im statischen
  Teil \ra die stärkste Schraube \\
$C_{load}$ & Fläche, Anzahl Transistoren, kleinere Strukturen \\
$\alpha,\ f$ & Taktfrequenz und Schaltfrequenz reduzieren, Glitches
  vermeiden, schlafen. $\alpha$ ist der Anteil der Gatter, die pro Takt
  wirklich umschalten \\
$t_{SC}$ & steile Flanken \ra kürzerer Kurzschluss \\
$I_{leak}$ & neue Halbleiter, Strukturen \emph{nicht} zu klein, gut
  isolierende Materialien, Temperatur tief halten \\
\end{ztab}
\begin{bsp}[Was die Quadratur bringt]
$V_{dd}$ von 3.3\,V auf 1.8\,V: der dynamische Anteil fällt auf
$(1.8/3.3)^2 \approx 30\,\%$, der statische auf $55\,\%$. Halber Takt
dagegen halbiert nur den dynamischen Anteil. Deshalb ist die erste
Messaufgabe genau diese: Einfluss von $V_{dd}$ beim LPC804.
\end{bsp}
